% !TeX root = ../main.tex
\chapter{Données brutes et scripts de traitement}
\label{app:graphes}

Les séries temporelles brutes d'acquisition enregistrées lors des travaux pratiques sont archivées au format LabVIEW (\texttt{.lvm}) dans le répertoire \texttt{Resultat/} du projet :
\begin{itemize}[leftmargin=1.2cm]
    \item \texttt{Resultat/Normal.lvm} : série temporelle brute de la température en flux d'air chaud cadencée à \qty{20}{\milli\second} ($N = 520$ points) ;
    \item \texttt{Resultat/Oversampling.lvm} : série temporelle brute de la température en flux d'air chaud avec suréchantillonnage moyenné par paquets de 100 points ($N = 375$ points) ;
    \item \texttt{Resultat/4w\_sauttemp.lvm} : réponse indicielle de la sonde RTD Pt100 lors du plongeon dans l'eau chaude ($N = 114$ points après calage du zéro initial).
\end{itemize}

Les scripts d'analyse numérique et de génération des figures vectorielles haute résolution respectant la charte graphique de la check-list TechMes sont disponibles dans le répertoire de code :
\begin{itemize}[leftmargin=1.2cm]
    \item \texttt{NTC.py} (dans \texttt{script/}) : chargement des fichiers LVM, calculs statistiques et tracé comparatif généré (figure~\ref{fig:diff_over_normal}) ;
    \item \texttt{generate\_graph\_pt100.py} (dans \texttt{scripts/}) : calage temporel, ajustement non linéaire du premier ordre et tracé indiciel complet (figure~\ref{fig:reponse_indicielle_pt100}).
\end{itemize}
